\documentclass[11pt,letterpaper]{article}
\usepackage[utf8]{inputenc}
\usepackage{geometry}
\usepackage{setspace}

\geometry{left=1.0in,right=1.0in,top=0.85in,bottom=0.85in}
\setstretch{1.25}

\setlength{\parindent}{0pt}
\setlength{\parskip}{8pt plus 1pt minus 1pt}

\begin{document}

\thispagestyle{empty}

\begin{center}
    \large\textbf{IN THE UNITED STATES PATENT AND TRADEMARK OFFICE} \\
\end{center}

\vspace{0.2cm}
\noindent\begin{tabular}{p{2.3in}p{4.0in}}
\textbf{In re Patent Application of:} & Reinaldo Javier Aguiar Marcano \\
\textbf{Application Number:} & 19/560,062 \\
\textbf{Confirmation Number:} & 9071 \\
\textbf{Filing Date:} & March 8, 2026 \\
\textbf{Group Art Unit:} & Pre-Exam / TC 2100 \\
\textbf{Examiner:} & Andrea Bennett \\
\textbf{Title of Invention:} & PORTABLE DATA CARTRIDGE WITH OPTICAL INTERFACE \\
\end{tabular}

\vspace{0.6cm}
\begin{center}
    \textbf{\large REPLACEMENT ABSTRACT (37 CFR \S{} 1.72(b))} \\
    \textbf{(COMMENCING ON A SEPARATE SHEET)}
\end{center}

\vspace{0.4cm}

\begin{center}
    \textbf{ABSTRACT}
\end{center}

A portable, self-contained, high-capacity data storage cartridge designed for the rapid physical transport of large datasets. The apparatus comprises a ruggedized housing containing a high-density, non-volatile storage medium, an integrated power source (e.g., a rechargeable battery), and at least one massively parallel optical interface for data ingest and offload. The optical interface, consistent with the technology of the parent application, allows the cartridge to be loaded or unloaded with petabyte-scale data in minutes, rather than the hours or days required with conventional interfaces. The cartridge is the foundational hardware for a new type of physical data logistics network, making the physical transfer of data a viable and often faster alternative to network transfer for bulk datasets.

\end{document}
